\magnification=\magstep1
\hsize=16truecm
\vsize=24truecm
\parindent=0pt
\baselineskip=14pt
\raggedright
\tolerance=10000
\hbadness=10000
\font\tenrm="Liberation Serif:mapping=tex-text" at 10pt
\font\tenbf="Liberation Serif/B:mapping=tex-text" at 10pt
\font\tenit="Liberation Serif/I:mapping=tex-text" at 10pt
\rm
\def\boxnum#1{\hbox{\vrule\vbox{\hrule\hbox{\strut\kern5pt{\bf#1}\kern5pt}\hrule}\vrule}}
\def\problem#1#2{\vskip0pt plus 30pt\penalty-400\bigskip
  \line{\boxnum{#1}\quad #2\hfil}\nobreak\smallskip}
\def\st#1{\medskip\noindent{\bf #1}\enspace\ignorespaces}
\def\ans#1{\medskip\noindent\hbox{\vrule\vbox{\hrule\hbox{\strut\kern5pt{\bf Ответ:}\enspace$\displaystyle #1$\kern5pt}\hrule}\vrule}\par}

\centerline{\bf Семинар 11. Линейные неоднородные уравнения}
\centerline{\bf второго порядка. Метод вариации постоянных}
\bigskip
{\bf Памятка: зачем ещё один метод?} Для полиномов, экспонент, синусов
и косинусов мы умеем выбирать пробное решение. Но, например, для
$1/\cos x$ такая таблица не даёт ответа. Вернёмся к идее, знакомой по
ЛНДУ первого порядка: заменить постоянную интегрирования функцией.
Теперь постоянных две, поэтому нужны две функции.

Рассмотрим нормированное уравнение
$$y''+p(x)y'+q(x)y=f(x).$$
Если перед $y''$ стоит $a(x)\ne0$, сначала делим {\it всё} уравнение на
$a(x)$, включая правую часть. Работаем на интервале, где коэффициенты и
$f$ непрерывны. Пусть уже известно общее решение ЛОДУ
$$y_h=C_1y_1(x)+C_2y_2(x).$$
Пробуем искать частное решение в той же форме, но с меняющимися множителями:
$$y_p=C_1(x)y_1(x)+C_2(x)y_2(x).$$
При постоянных множителях левая часть равна нулю; изменения множителей
должны создать нужную правую часть.

\st{Откуда дополнительное условие?} Одну функцию $y_p$ мы представили
через две неизвестные функции. Такое представление избыточно: изменение
первого слагаемого можно компенсировать вторым. Используем эту свободу,
чтобы сделать производную похожей на производную однородного решения.
По правилу произведения
$$y_p'=C_1'y_1+C_2'y_2+C_1y_1'+C_2y_2'.$$
Выбираем функции так, чтобы
$$C_1'y_1+C_2'y_2=0.$$
Это не новое условие на само уравнение или его начальные данные, а удобный
выбор представления решения. Его выполнимость подтвердится разрешимостью
системы ниже. Тогда
$$y_p'=C_1y_1'+C_2y_2',\qquad
y_p''=C_1'y_1'+C_2'y_2'+C_1y_1''+C_2y_2''.$$
Подставляем в ЛНДУ и группируем:
$$C_1'y_1'+C_2'y_2'+C_1(y_1''+py_1'+qy_1)
 +C_2(y_2''+py_2'+qy_2)=f.$$
Скобки равны нулю, так как $y_1,y_2$ решают ЛОДУ. Получаем систему
$$\cases{y_1C_1'+y_2C_2'=0,\cr
             y_1'C_1'+y_2'C_2'=f.\cr}$$

\st{Как решить систему?} Её определитель
$$W=y_1y_2'-y_2y_1'$$
называется вронскианом. В наших примерах проверяем $W\ne0$ на выбранном
интервале. Умножим первое уравнение на $y_2'$, второе на $y_2$ и вычтем:
$WC_1'=-y_2f$. Аналогично, второе умножим на $y_1$, первое на $y_1'$:
$WC_2'=y_1f$. Следовательно
$$C_1'=-{y_2f\over W},\qquad C_2'={y_1f\over W}.$$
Это обычная алгебраическая система относительно производных, а не новая
система дифференциальных уравнений, которую надо решать особым методом.

\st{Интегрирование и ответ.}
$$C_1(x)=-\int {y_2(x)f(x)\over W(x)}\,dx,\qquad
C_2(x)=\int {y_1(x)f(x)\over W(x)}\,dx.$$
Для одного частного решения постоянные этих интегралов берём равными нулю.
Если их сохранить, они дадут уже известные слагаемые $y_h$.
Итак, $y=y_h+C_1(x)y_1+C_2(x)y_2$. Не путайте функции $C_1(x),C_2(x)$
в частном решении с произвольными {\it числами} $C_1,C_2$ в общем ответе.
Начальные условия накладываем на полное решение, не только на $y_h$.

\st{Алгоритм.} Нормировать уравнение; найти или взять из условия $y_1,y_2$;
проверить $W\ne0$; записать два уравнения на $C_1',C_2'$; решить их;
интегрировать; собрать $y_p$, затем $y_h+y_p$; указать интервал и проверить
подстановкой. Формулы выше -- итог вывода, а не подстановка без объяснения.

\vfil\eject
\centerline{\bf Классная работа: условия и подробные решения}
\bigskip
\abovedisplayskip=7pt plus 2pt minus 2pt
\belowdisplayskip=7pt plus 2pt minus 2pt

\problem{1}{$y''+y=1/\cos x,\qquad |x|<\pi/2$}
\st{Шаг 1.} Уравнение уже нормировано: коэффициент при $y''$ равен 1.
Решаем однородное уравнение:
$$y''+y=0,\qquad\lambda^2+1=0\quad\Rightarrow\quad\lambda=\pm i.$$
$$y_1=\cos x,\quad y_2=\sin x,\qquad y_h=C_1\cos x+C_2\sin x.$$
Находим производные и вронскиан:
$$y_1'=-\sin x,\quad y_2'=\cos x,\qquad
W=\cos x\cdot\cos x-\sin x\cdot(-\sin x)=1\ne0.$$
\st{Шаг 2.} Заменяем постоянные функциями:
$$y_p=C_1(x)\cos x+C_2(x)\sin x.$$
По правилу произведения
$$y_p'=C_1'\cos x+C_2'\sin x-C_1\sin x+C_2\cos x.$$
Накладываем условие $C_1'\cos x+C_2'\sin x=0$. Тогда
$$y_p'=-C_1\sin x+C_2\cos x,$$
$$y_p''=-C_1'\sin x+C_2'\cos x-C_1\cos x-C_2\sin x.$$
При сложении с $y_p$ последние два слагаемых сокращаются:
$$y_p''+y_p=-C_1'\sin x+C_2'\cos x={1\over\cos x}.$$
\st{Шаг 3.} Получили систему на производные неизвестных функций:
$$\cases{C_1'\cos x+C_2'\sin x=0,\cr
          -C_1'\sin x+C_2'\cos x=1/\cos x.\cr}$$
Первое уравнение умножаем на $\cos x$, второе на $-\sin x$ и складываем:
$$C_1'(\cos^2x+\sin^2x)=-{\sin x\over\cos x}
\quad\Rightarrow\quad C_1'=-\tan x.$$
Теперь первое умножаем на $\sin x$, второе на $\cos x$ и складываем:
$$C_2'(\sin^2x+\cos^2x)=1\quad\Rightarrow\quad C_2'=1.$$
\vfil\eject
\line{\boxnum{1}\quad Продолжение решения\hfil}
\smallskip
\st{Шаг 4.} Интегрируем; для частного решения постоянные берём нулевыми.
При $|x|<\pi/2$ имеем $\cos x>0$, поэтому
$$C_1(x)=-\int{\sin x\over\cos x}\,dx
=\int{d(\cos x)\over\cos x}=\ln\cos x,\qquad C_2(x)=\int dx=x.$$
\st{Шаг 5.} Подставляем функции на место множителей:
$$y_p=(\ln\cos x)\cos x+x\sin x,\qquad y=y_h+y_p.$$
\st{Проверка.} Дифференцируем найденное частное решение:
$$y_p'=-\sin x\ln\cos x-\sin x+\sin x+x\cos x
=-\sin x\ln\cos x+x\cos x,$$
$$y_p''=-\cos x\ln\cos x+\sin x\tan x+\cos x-x\sin x.$$
Складываем с $y_p$ и сокращаем логарифмические слагаемые и $x\sin x$:
$$y_p''+y_p={\sin^2x\over\cos x}+\cos x
={\sin^2x+\cos^2x\over\cos x}={1\over\cos x}.$$
\ans{y=C_1\cos x+C_2\sin x+\cos x\ln\cos x+x\sin x}

\vfil\eject
\problem{2}{$y''+y=\tan x,\qquad |x|<\pi/2$}
\st{Шаг 1.} Решаем однородное уравнение:
$$\lambda^2+1=0\quad\Rightarrow\quad\lambda=\pm i,\qquad
y_h=C_1\cos x+C_2\sin x.$$
Берём $y_1=\cos x$, $y_2=\sin x$; тогда
$$W=\cos x\cdot\cos x-\sin x\cdot(-\sin x)=1.$$
\st{Шаг 2.} Ищем $y_p=C_1(x)\cos x+C_2(x)\sin x$ и накладываем
условие $C_1'\cos x+C_2'\sin x=0$. При этом
$$y_p'=-C_1\sin x+C_2\cos x,$$
$$y_p''=-C_1'\sin x+C_2'\cos x-C_1\cos x-C_2\sin x.$$
Подставляем в $y_p''+y_p=\tan x$: слагаемые с $C_1,C_2$ сокращаются.
Остаётся система
$$\cases{C_1'\cos x+C_2'\sin x=0,\cr
          -C_1'\sin x+C_2'\cos x=\tan x.\cr}$$
\st{Шаг 3.} Для исключения $C_2'$ складываем первое уравнение,
умноженное на $\cos x$, и второе, умноженное на $-\sin x$:
$$C_1'=-\sin x\tan x=-{\sin^2x\over\cos x}
=-{1-\cos^2x\over\cos x}=\cos x-{1\over\cos x}.$$
Для исключения $C_1'$ складываем первое, умноженное на $\sin x$,
и второе, умноженное на $\cos x$:
$$C_2'=\cos x\tan x=\sin x.$$
\vfil\eject
\line{\boxnum{2}\quad Продолжение решения\hfil}
\smallskip
\st{Шаг 4.} Разбираем нетабличный на первый взгляд интеграл.
Умножаем числитель и знаменатель на $\cos x$ и делаем замену $t=\sin x$:
$$\int{dx\over\cos x}=\int{\cos x\,dx\over1-\sin^2x}
=\int{dt\over1-t^2}.$$
Разложение на простые дроби даёт
$$\int{dt\over1-t^2}={1\over2}\int\left({1\over1+t}+{1\over1-t}\right)dt
={1\over2}\ln{1+t\over1-t}.$$
На выбранном интервале $-1<\sin x<1$, $\cos x>0$. Поэтому
$$H(x):={1\over2}\ln{1+\sin x\over1-\sin x}
=\ln{1+\sin x\over\cos x}=\ln(\sec x+\tan x),$$
где $\sec x=1/\cos x$. Итак, $H'=1/\cos x$ и
$$C_1(x)=\int\cos x\,dx-\int{dx\over\cos x}=\sin x-H(x),$$
$$C_2(x)=\int\sin x\,dx=-\cos x.$$
\st{Шаг 5.} Собираем частное решение, явно сокращая слагаемые:
$$y_p=(\sin x-H)\cos x+(-\cos x)\sin x=-H\cos x.$$
Полное решение $y=y_h+y_p$.
\st{Проверка.} Используем $H'=1/\cos x$:
$$y_p'=\sin x\,H-\cos x\,{1\over\cos x}=\sin x\,H-1,$$
$$y_p''=\cos x\,H+{\sin x\over\cos x},\qquad
y_p''+y_p=\cos x\,H+\tan x-\cos x\,H=\tan x.$$
\ans{y=C_1\cos x+C_2\sin x-\cos x\ln(\sec x+\tan x)}

\vfil\eject
\problem{3}{$y''=1/(1+x^2),\qquad x\in\mathop{\rm R}$}
\st{Шаг 1.} Однородное уравнение $y''=0$ даёт
$$\lambda^2=0,\qquad\lambda_1=\lambda_2=0,\qquad
y_h=(C_1+C_2x)e^{0x}=C_1+C_2x.$$
Выбираем $y_1=1$, $y_2=x$. Проверяем независимость:
$$y_1'=0,\quad y_2'=1,\qquad W=1\cdot1-x\cdot0=1.$$
Здесь можно дважды интегрировать исходное уравнение; решим его вариацией
постоянных, чтобы сопоставить методы.
\st{Шаг 2.} Заменяем постоянные функциями:
$$y_p=C_1(x)+xC_2(x),\qquad y_p'=C_1'+xC_2'+C_2.$$
Накладываем условие $C_1'+xC_2'=0$. Тогда $y_p'=C_2$, $y_p''=C_2'$.
Подстановка в исходное уравнение даёт систему
$$\cases{C_1'+xC_2'=0,\cr C_2'=1/(1+x^2).\cr}$$
\st{Шаг 3.} Подставляем второе уравнение в первое:
$$C_1'+{x\over1+x^2}=0\quad\Rightarrow\quad
C_1'=-{x\over1+x^2}.$$
Интегрируем первое выражение заменой $z=1+x^2$, $dz=2x\,dx$:
$$C_1(x)=-\int{x\,dx\over1+x^2}=-{1\over2}\int{dz\over z}
=-{1\over2}\ln z=-{1\over2}\ln(1+x^2).$$
Для второго используем табличный интеграл:
$$C_2(x)=\int{dx\over1+x^2}=\arctan x.$$
\st{Шаг 4.} Возвращаемся к $y_p$ и добавляем $y_h$:
$$y_p=-{1\over2}\ln(1+x^2)+x\arctan x,\qquad y=y_h+y_p.$$
\st{Проверка.} Дифференцируем оба слагаемых:
$$y_p'=-{1\over2}\,{2x\over1+x^2}+\arctan x+{x\over1+x^2}
=\arctan x,$$
$$y_p''={1\over1+x^2},\qquad y_h''=(C_1+C_2x)''=0.$$
Так как $1+x^2>0$ для любого действительного $x$, решение определено
на всей действительной оси.
\ans{y=C_1+C_2x-{1\over2}\ln(1+x^2)+x\arctan x}

\vfil\eject
\problem{4}{$y''-y=e^x/(1+e^x),\qquad x\in\mathop{\rm R}$}
\st{Шаг 1.} Решаем однородное уравнение $y''-y=0$:
$$\lambda^2-1=(\lambda-1)(\lambda+1)=0
\quad\Rightarrow\quad\lambda_1=1,\quad\lambda_2=-1.$$
$$y_1=e^x,\quad y_2=e^{-x},\qquad y_h=C_1e^x+C_2e^{-x}.$$
Вычисляем вронскиан:
$$W=e^x(-e^{-x})-e^{-x}e^x=-1-1=-2\ne0.$$
\st{Шаг 2.} Ищем $y_p=C_1(x)e^x+C_2(x)e^{-x}$.
Условие $C_1'e^x+C_2'e^{-x}=0$ сокращает первую производную до
$$y_p'=C_1e^x-C_2e^{-x}.$$
Тогда
$$y_p''=C_1'e^x-C_2'e^{-x}+C_1e^x+C_2e^{-x},$$
$$y_p''-y_p=C_1'e^x-C_2'e^{-x}={e^x\over1+e^x}.$$
Получаем систему
$$\cases{C_1'e^x+C_2'e^{-x}=0,\cr
          C_1'e^x-C_2'e^{-x}=e^x/(1+e^x).\cr}$$
\st{Шаг 3.} Складываем уравнения, а затем вычитаем второе из первого:
$$2C_1'e^x={e^x\over1+e^x}
\quad\Rightarrow\quad C_1'={1\over2(1+e^x)},$$
$$2C_2'e^{-x}=-{e^x\over1+e^x}
\quad\Rightarrow\quad C_2'=-{e^{2x}\over2(1+e^x)}.$$
\vfil\eject
\line{\boxnum{4}\quad Продолжение решения\hfil}
\smallskip
\st{Шаг 4.} Для первого интеграла выделяем производную логарифма:
$$C_1(x)={1\over2}\int{dx\over1+e^x}
={1\over2}\int\left(1-{e^x\over1+e^x}\right)dx$$
$$={1\over2}\left(x-\int{d(1+e^x)\over1+e^x}\right)
={x-\ln(1+e^x)\over2}.$$
Во втором интеграле делаем замену $z=e^x$, $dz=e^x dx$, $dx=dz/z$:
$$C_2(x)=-{1\over2}\int{e^{2x}\over1+e^x}\,dx
=-{1\over2}\int{z^2\over1+z}\,{dz\over z}.$$
Делим $z$ на $1+z$: ${z\over1+z}=1-{1\over1+z}$. Поэтому
$$C_2(x)=-{1\over2}\int{z\over1+z}\,dz
 =-{1\over2}\bigl[z-\ln(1+z)\bigr]={-e^x+\ln(1+e^x)\over2}.$$
\st{Шаг 5.} Для краткости обозначим $H=\ln(1+e^x)$. Подставляем:
$$y_p={x-H\over2}\,e^x+{-e^x+H\over2}\,e^{-x}
={e^x(x-H)-1+e^{-x}H\over2}.$$
Добавляем $y_h$. Здесь $1+e^x>0$ при всех $x$, ограничений нет.
\st{Проверка.} Сначала проверяем производные найденных функций:
$$C_1'={1\over2}\left(1-{e^x\over1+e^x}\right)
={1\over2(1+e^x)},$$
$$C_2'={1\over2}\left(-e^x+{e^x\over1+e^x}\right)
=-{e^{2x}\over2(1+e^x)}.$$
Их подстановка в дополнительное условие даёт
$$C_1'e^x+C_2'e^{-x}={e^x\over2(1+e^x)}-{e^x\over2(1+e^x)}=0.$$
Следовательно, формулы производных из шага 2 применимы, и
$$y_p''-y_p=C_1'e^x-C_2'e^{-x}
={e^x\over2(1+e^x)}+{e^x\over2(1+e^x)}={e^x\over1+e^x}.$$
\ans{y=C_1e^x+C_2e^{-x}+{e^x(x-H)+e^{-x}H-1\over2}}

\vfil\eject
\problem{5}{$x^2y''-2xy'+2y=x^2,\qquad x>0$}
\st{Шаг 1.} Делим всё уравнение на $x^2\ne0$:
$$y''-{2\over x}y'+{2\over x^2}y=1.$$
Теперь $p(x)=-2/x$, $q(x)=2/x^2$, $f(x)=1$.
В системе вариации постоянных справа будет $1$, а не $x^2$.
\st{Шаг 2.} Для однородной части даны $y_1=x$, $y_2=x^2$.
Проверяем их в исходном однородном уравнении:
$$y_1'=1,\quad y_1''=0,\qquad x^2\cdot0-2x\cdot1+2x=0,$$
$$y_2'=2x,\quad y_2''=2,\qquad x^2\cdot2-2x\cdot2x+2x^2=0.$$
Вронскиан $W=x\cdot2x-x^2\cdot1=x^2\ne0$ при $x>0$.
Поэтому $y_h=C_1x+C_2x^2$.
\st{Шаг 3.} Ищем $y_p=xC_1(x)+x^2C_2(x)$. Тогда
$$y_p'=xC_1'+x^2C_2'+C_1+2xC_2.$$
Выбираем $xC_1'+x^2C_2'=0$, после чего
$$y_p'=C_1+2xC_2,\qquad y_p''=C_1'+2xC_2'+2C_2.$$
Подставляем в нормированное уравнение:
$$C_1'+2xC_2'+2C_2-{2\over x}(C_1+2xC_2)
+{2\over x^2}(xC_1+x^2C_2)=1,$$
$$C_1'+2xC_2'+(2-4+2)C_2+\left(-{2\over x}+{2\over x}\right)C_1=1.$$
Получаем систему
$$\cases{xC_1'+x^2C_2'=0,\cr C_1'+2xC_2'=1.\cr}$$
\vfil\eject
\line{\boxnum{5}\quad Продолжение решения\hfil}
\smallskip
\st{Шаг 4.} Делим первое уравнение на $x>0$:
$$C_1'+xC_2'=0.$$
Вычитаем его из второго:
$$xC_2'=1\quad\Rightarrow\quad C_2'={1\over x},\qquad
C_1'=-xC_2'=-1.$$
Интегрируем:
$$C_1(x)=-\int dx=-x,\qquad C_2(x)=\int{dx\over x}=\ln x.$$
\st{Шаг 5.} Подставляем в $y_p$, затем добавляем $y_h$:
$$y_p=x(-x)+x^2\ln x=-x^2+x^2\ln x,$$
$$y=C_1x+C_2x^2-x^2+x^2\ln x=C_1x+(C_2-1)x^2+x^2\ln x.$$
Число $C_2-1$ произвольно; обозначаем его снова $C_2$.
\st{Проверка.} Достаточно проверить частное слагаемое $v=x^2\ln x$:
$$v'=2x\ln x+x,\qquad v''=2\ln x+2+1=2\ln x+3.$$
В исходном, а не нормированном уравнении получаем
$$x^2v''-2xv'+2v=x^2(2\ln x+3)-2x(2x\ln x+x)+2x^2\ln x$$
$$=(2-4+2)x^2\ln x+(3-2)x^2=x^2.$$
\ans{y=C_1x+C_2x^2+x^2\ln x,\quad x>0}

\vfil\eject
\problem{6}{$y''+y=1/\cos x,\quad y(0)=0,\ y'(0)=0$}
\st{Шаг 1.} Для однородного уравнения
$$\lambda^2+1=0\quad\Rightarrow\quad\lambda=\pm i,\qquad
y_h=C_1\cos x+C_2\sin x.$$
Берём $y_1=\cos x$, $y_2=\sin x$, $W=\cos^2x+\sin^2x=1$.
Выбираем интервал $|x|<\pi/2$: он содержит начальную точку $0$,
а $\cos x$ на нём положителен.
\st{Шаг 2.} Ищем $y_p=C_1(x)\cos x+C_2(x)\sin x$.
С дополнительным условием первая производная имеет вид
$$y_p'=-C_1(x)\sin x+C_2(x)\cos x.$$
Подстановка второй производной в $y_p''+y_p=1/\cos x$ даёт
$$\cases{C_1'\cos x+C_2'\sin x=0,\cr
          -C_1'\sin x+C_2'\cos x=1/\cos x.\cr}$$
Умножая уравнения соответственно на $\cos x,-\sin x$ и складывая,
а затем на $\sin x,\cos x$ и складывая, получаем
$$C_1'(\cos^2x+\sin^2x)=-{\sin x\over\cos x},\qquad
C_2'(\sin^2x+\cos^2x)=1.$$
Следовательно,
$$C_1'=-\tan x,\quad C_2'=1,\qquad
C_1(x)=\ln\cos x,\quad C_2(x)=x.$$
\st{Шаг 3.} Собираем частное и полное решения:
$$y_p=\cos x\ln\cos x+x\sin x,$$
$$y=C_1\cos x+C_2\sin x+\cos x\ln\cos x+x\sin x.$$
Вычисляем производную полного решения:
$$y'=-C_1\sin x+C_2\cos x-\sin x\ln\cos x+x\cos x.$$
Здесь слагаемые $-\sin x$ от производной $\cos x\ln\cos x$
и $+\sin x$ от производной $x\sin x$ сократились.
\vfil\eject
\line{\boxnum{6}\quad Продолжение решения\hfil}
\smallskip
\st{Шаг 4.} Подставляем $x=0$, учитывая $\ln\cos0=\ln1=0$:
$$y(0)=C_1\cdot1+C_2\cdot0+1\cdot0+0\cdot0=C_1=0,$$
$$y'(0)=-C_1\cdot0+C_2\cdot1-0\cdot0+0\cdot1=C_2=0.$$
Обе произвольные постоянные равны нулю: $y=y_p$.
\st{Проверка.} Для найденной функции
$$y'=-\sin x\ln\cos x+x\cos x,$$
$$y''=-\cos x\ln\cos x+{\sin^2x\over\cos x}+\cos x-x\sin x,$$
$$y''+y={\sin^2x+\cos^2x\over\cos x}={1\over\cos x},$$
$$y(0)=1\cdot\ln1+0=0,\qquad y'(0)=-0\cdot\ln1+0=0.$$
Через $x=\pm\pi/2$ решение задачи продолжить нельзя:
правая часть исходного уравнения там не определена.
\ans{y=\cos x\ln\cos x+x\sin x,\quad |x|<\pi/2}

\vfil\eject
\centerline{\bf Самостоятельная работа}
\bigskip
Примените вариацию постоянных, выпишите систему на $C_1'(x),C_2'(x)$.
В последней задаче однородная часть дана; предварительно нормируйте
уравнение. Проверьте решение на указанном интервале.
\problem{1}{$y''+y=2/\cos x,\quad |x|<\pi/2$}
\problem{2}{$y''+y=1/\sin x,\quad 0<x<\pi$}
\problem{3}{$y''=1/x,\quad x>0$}
\problem{4}{$y''=\ln x,\quad x>0$}
\problem{5}{$y''+y=\tan x,\quad y(0)=1,\ y'(0)=0,\quad |x|<\pi/2$}
\problem{6}{$x^2y''-2xy'+2y=x^3,\quad x>0$}
Для задачи 6: $y_h=C_1x+C_2x^2$.
\bye
